% ======================================================================
\section{Discussion and reflection}\label{sec:discussion}
% ======================================================================
\paragraph{Which algorithm, where, and why} Grounded in \cref{tab:main,fig:convergence,fig:stn,fig:extended}:
\emph{(i)} On small and medium instances TS and MA are equivalent in final quality (Wilcoxon $p=0.22$ over 73
instances). MA finds its final solution sooner on the larger A instances (0.02\,s vs.\ 2.5\,s on A-n60-k9, 0.8\,s vs.\
27.6\,s on A-n80-k10), whereas TS is faster and more reliable on the team instance (median 4.1\,s vs.\ 15.1\,s to reach
0.1\,\%; 9 vs.\ 7 runs at the best value). \emph{(ii)} TS starts best: it begins from the Clarke--Wright plan reduced to $K$ routes,
and where that start violates capacity (A-n45-k6) its walk through infeasible space (49\,\% of all iterations,
\cref{app:extra}) restores feasibility within about ten iterations while it already shortens the routes. \emph{(iii)} On the tightest instances TS is the only safe choice: with a load factor of 0.992
(P-n55-k15) neither GA nor MA ever produced a plan within the fleet, because their decoders and local search stay in
feasible space, while TS crosses infeasible space and reached the optimum in every run. \emph{(iv)} The pure GA is not
competitive at equal time, and its failure mode is structural rather than a tuning issue (\cref{fig:stn}): permutation
operators rarely realise the partition-level change the CVRP needs. \emph{(v)} \emph{Size changes the ranking.} Up to $n\approx200$ TS and MA are on par, but from $n=250$ to $1000$ MA is better on every X instance (mean gap 1.12\,\% vs.\ 2.04\,\% for $n>300$): best-improvement TS pays $\Theta(nk)$ evaluations per applied move (\cref{sec:complexity}), which is affordable for 100 clinics and not for 1000 within one minute.


\paragraph{How the Bosphorus shaped the routes} Less than we expected on our own instance, and exactly as the theory
says. With the laboratory on the strait and the European clusters far to the west, the pure-distance optimum already
uses the minimum number of crossings, so the penalty adds $20P$ to the bill and changes nothing else
(\cref{cor:invariant}); all 100 sweep runs agree. On real geography, where clinics sit close to both shores, the penalty
does reshape the plan: above $P^\ast=1.12$\,km one crossing disappears, exactly where \cref{prop:cross}(d) predicts.
In both cases the number of crossing routes equals the bin-packing bound ($10$ and $6$) and no route ever zig-zags.
Fully side-pure routing would need $\BP(q|_E)+\BP(q|_A)=10+8=18=K$ drones on the team instance --- possible only
because of the spare vehicle --- but the optimum prefers 17 drones with one mixed route: the theory pins down the
\emph{number} of crossings, while the optimiser still trades a vehicle against a few kilometres.

\paragraph{Theory vs.\ practice} The best-rotation tour-partitioning heuristic of \cref{sec:approx} carries a
\emph{proven} factor of $3.96$; on our six main instances its actual ratios are only 1.09--1.21, computed in
5--30\,ms --- but it ignores the fleet limit and uses one or two drones too many on five of them. TS and MA carry
\emph{no} guarantee, yet 112 of their 120 runs on the six main instances ended within 0.1\,\% of the best known value. Guarantees are
worst-case statements over all metric instances, including adversarial ones that a dispatcher never meets;
metaheuristics are tuned to the distribution of instances that actually occur and can absorb side constraints
(fleet size, a strait, later time windows) without a new proof. Industry accepts the trade-off for three reasons we
reproduced here: the empirical quality is far better than the bound, the running time is under control (anytime), and
quality can be \emph{certified a posteriori} --- with a lower bound (here: the radial bound and the bin-packing bound
of \cref{prop:cross}) or, for small cases, an exact solver. Theory still matters: it told us that feasibility itself is
hard (so we designed the penalty), that the strait keeps the metric (so the approximation results and
Clarke--Wright's merge logic remain valid), and where the hidden bin packing sits.

\paragraph{Limitations} Instances with up to 100 clinics (1000 for the X study); deterministic demands; a constant
crossing penalty (a corridor-based penalty via the bridges would be the shortest-path closure of \cref{thm:metric}'s
construction and remains metric); no battery-range or time-window constraints; one machine.

\subsection*{Personal reflection}
\textbf{What surprised us most.} That the recommended Clarke--Wright seed can trap the GA: one strong individual took
over the population and ten different random seeds produced the identical result 827 on A-n32-k5, while the unseeded
GA reached the optimum in eight of ten runs --- yet on the five other instances the same seed helps. And that the Bosphorus,
the centre of our story, does not change a single route of our team instance --- we expected a gradual ``phase
diagram'', the corollary explains in one line why there is none, and real geography then produced exactly the
transition we had looked for. Second, that the
penalty update behaves like a thermostat: TS spends 49\,\% of its iterations in feasible space. In hindsight this is
forced by the rule itself --- $\log\lambda$ moves by $\pm\log(1+\delta)$ every iteration, so a bounded $\lambda$ needs as
many feasible as infeasible iterations, whatever $\delta$.
\textbf{The hardest bug} was not in a solver but in an analysis script, and it nearly produced a wrong conclusion. The
first search-trajectory plot showed every TS run as just two points, as if TS did nothing after its first feasible
solution. The plotting code coloured each partition with the \emph{best} gap ever observed for it instead of the gap
\emph{at the time of the visit}, so all sequencing improvements inside one partition collapsed into a single point. We
found it because the picture contradicted the run logs (TS improves from 989 to 944 on A-n45-k6); the fix exposed the
vertical ``sequencing drop'' of \cref{fig:stn}b, which became one of our findings. In the solver itself, the subtlest
issue was the GA's sorted population: after a replacement the newcomer was bubbled left, but the right-bubbling loop
restarted from the old index, which could silently break elitism; a code review of the replacement step caught it
before any experiment was run.
\textbf{With more time} we would (1) combine the two winners --- a hybrid genetic search with TS as the education
operator and diversity-aware survival \citep{vidal2012hybrid,vidal2022hybrid}, which our STN analysis suggests would fix
the GA's takeover; (2) replace the constant $P$ by a corridor model through the three bridges and add battery range
(payload-dependent energy as in \citep{dorling2017vehicle}) and time windows for blood samples; (3) learn instance-specific
neighbourhood sizes, in the spirit of recent LLM-designed heuristics \citep{xie2026enhancing}; and (4) feed risk maps
from drone-vision pipelines such as \citep{alpay2026pattern} into time-dependent crossing penalties.
